% --- Archon physics formalization source begin ---
% archon:physics
% archon:covers PhyXMiniProblems/problem_phyx_mini_0821.lean
% archon:source-report reports/phyx_mini/problem_phyx_mini_0821.source.json

\chapter{Physics problem phyx\_mini\_0821}
\label{ch:PhyXMiniProblems_problem_phyx_mini_0821}

\paragraph{Problem source.}
\#\# Physical scenario

We wrap a light, nonstretching cable around a solid cylinder with mass \textbackslash{}( M \textbackslash{}) and radius \textbackslash{}( R \textbackslash{}). The cylinder rotates with negligible friction about a stationary horizontal axis. We tie the free end of the cable to a block of mass \textbackslash{}( m \textbackslash{}) and release the block from rest at a distance \textbackslash{}( h \textbackslash{}) above the floor as shown in figure. As the block falls, the cable unwinds without stretching or slipping.

\#\# Question

Find the angular speed of the cylinder as the block strikes the floor.

\#\# Answer choices

- A: A: \textbackslash{}(2v/R\textbackslash{})

- B: B: \textbackslash{}(v/R\textbackslash{})

- C: C: \textbackslash{}(v/2R\textbackslash{})

- D: D: \textbackslash{}(2v/3R\textbackslash{})

\#\# Figure caption (auxiliary; use the image as primary evidence)

The image depicts a physics diagram involving a pulley system. Here's a detailed description:

- **Pulley**: At the top, there is a circular pulley labeled with a radius "R" and a mass "M."
- **String and Block**: A string is wound around the pulley and connected to a block. The block is labeled "m," indicating its mass.
- **Height**: Below the block, there is a measurement labeled "h," indicating the distance from the block to the ground surface.
- **Ground**: The ground is represented by a horizontal line with diagonal hatching below it, indicating the surface on which the block will land.

The system appears to illustrate a basic mechanics problem involving potential and kinetic energy, considering the masses of the pulley and the block, and the height from which the block can fall.

\#\# Dataset metadata

Rotational Motion; Physical Model Grounding Reasoning, Multi-Formula Reasoning

\paragraph{Recorded answer/context.}
B: B: \textbackslash{}(v/R\textbackslash{})

\paragraph{Figure/image path.}
/root/proposal\_for\_physic/science-mango/phyx\_mini\_run/phyx\_data/test\_image/821.png

\paragraph{Formalization target.}
create a compiling Lean file with sorry bodies at `PhyXMiniProblems/problem\_phyx\_mini\_0821.lean`. The Lean declarations must preserve the physical quantities, dimensions or dimensional roles, figure labels, governing-law hypotheses, and final relation expressed by this problem.
Use Mathlib/Physlib names found through LeanExplore where available. If a physics API is missing, introduce faithful local abstractions rather than scalar placeholder aliases.

\begin{theorem}[Physics formalization target]
\label{thm:physics:phyx_mini_0821:target}
\lean{PhyXMiniProblems.ProblemPhyXMini0821.angularSpeedAtImpact_eq_blockSpeed_div_radius}
\uses{def:physics:phyx-mini-0821:phyxminiproblems-problemphyxmini0821-lengthinmeters, def:physics:phyx-mini-0821:phyxminiproblems-problemphyxmini0821-speedinmeterspersecond, def:physics:phyx-mini-0821:phyxminiproblems-problemphyxmini0821-angularspeedinradianspersecond, def:physics:phyx-mini-0821:phyxminiproblems-problemphyxmini0821-cylindercablesetup, def:physics:phyx-mini-0821:phyxminiproblems-problemphyxmini0821-matchesprimaryfigure, def:physics:phyx-mini-0821:phyxminiproblems-problemphyxmini0821-matchesproblemdescription, def:physics:phyx-mini-0821:phyxminiproblems-problemphyxmini0821-hasphysicalparameters, def:physics:phyx-mini-0821:phyxminiproblems-problemphyxmini0821-satisfiescylindercablelaws}
The assigned autoformalize agent should translate this physics problem into Lean declarations in the covered file, with theorem and lemma proof bodies written as `by sorry`.
\end{theorem}
\begin{proof}
This is an autoformalization task, not a proof task. Produce statements that can later be proved by the physics prover without weakening the physical model.
\end{proof}

% --- Archon declaration topology begin ---
\section{Declaration topology}
\begin{definition}[acceleration Dimension]
  \label{def:physics:phyx-mini-0821:phyxminiproblems-problemphyxmini0821-accelerationdimension}
  \lean{PhyXMiniProblems.ProblemPhyXMini0821.accelerationDimension}
  The physical dimension L T⁻² of gravitational acceleration
\end{definition}
\begin{proof}
  This is the declared model definition of acceleration Dimension.
\end{proof}
\begin{definition}[moment Of Inertia Dimension]
  \label{def:physics:phyx-mini-0821:phyxminiproblems-problemphyxmini0821-momentofinertiadimension}
  \lean{PhyXMiniProblems.ProblemPhyXMini0821.momentOfInertiaDimension}
  The physical dimension M L² of an axial moment of inertia
\end{definition}
\begin{proof}
  This is the declared model definition of moment Of Inertia Dimension.
\end{proof}
\begin{definition}[Mass Quantity]
  \label{def:physics:phyx-mini-0821:phyxminiproblems-problemphyxmini0821-massquantity}
  \lean{PhyXMiniProblems.ProblemPhyXMini0821.MassQuantity}
  A nonnegative, unit-independent physical mass
\end{definition}
\begin{proof}
  This is the declared model definition of Mass Quantity.
\end{proof}
\begin{definition}[Length Quantity]
  \label{def:physics:phyx-mini-0821:phyxminiproblems-problemphyxmini0821-lengthquantity}
  \lean{PhyXMiniProblems.ProblemPhyXMini0821.LengthQuantity}
  A nonnegative, unit-independent physical length
\end{definition}
\begin{proof}
  This is the declared model definition of Length Quantity.
\end{proof}
\begin{definition}[Speed Quantity]
  \label{def:physics:phyx-mini-0821:phyxminiproblems-problemphyxmini0821-speedquantity}
  \lean{PhyXMiniProblems.ProblemPhyXMini0821.SpeedQuantity}
  A nonnegative linear-speed magnitude
\end{definition}
\begin{proof}
  This is the declared model definition of Speed Quantity.
\end{proof}
\begin{definition}[Angular Speed Quantity]
  \label{def:physics:phyx-mini-0821:phyxminiproblems-problemphyxmini0821-angularspeedquantity}
  \lean{PhyXMiniProblems.ProblemPhyXMini0821.AngularSpeedQuantity}
  A nonnegative angular-speed magnitude; radians are dimensionless
\end{definition}
\begin{proof}
  This is the declared model definition of Angular Speed Quantity.
\end{proof}
\begin{definition}[Acceleration Quantity]
  \label{def:physics:phyx-mini-0821:phyxminiproblems-problemphyxmini0821-accelerationquantity}
  \lean{PhyXMiniProblems.ProblemPhyXMini0821.AccelerationQuantity}
  \uses{def:physics:phyx-mini-0821:phyxminiproblems-problemphyxmini0821-accelerationdimension}
  A nonnegative gravitational-acceleration magnitude
\end{definition}
\begin{proof}
  This is the declared model definition of Acceleration Quantity.
\end{proof}
\begin{definition}[Moment Of Inertia Quantity]
  \label{def:physics:phyx-mini-0821:phyxminiproblems-problemphyxmini0821-momentofinertiaquantity}
  \lean{PhyXMiniProblems.ProblemPhyXMini0821.MomentOfInertiaQuantity}
  \uses{def:physics:phyx-mini-0821:phyxminiproblems-problemphyxmini0821-momentofinertiadimension}
  A nonnegative scalar moment of inertia about the cylinder axis
\end{definition}
\begin{proof}
  This is the declared model definition of Moment Of Inertia Quantity.
\end{proof}
\begin{definition}[Energy Quantity]
  \label{def:physics:phyx-mini-0821:phyxminiproblems-problemphyxmini0821-energyquantity}
  \lean{PhyXMiniProblems.ProblemPhyXMini0821.EnergyQuantity}
  A physical energy, with dimension M L² T⁻²
\end{definition}
\begin{proof}
  This is the declared model definition of Energy Quantity.
\end{proof}
\begin{definition}[nonnegative SIReadout]
  \label{def:physics:phyx-mini-0821:phyxminiproblems-problemphyxmini0821-nonnegativesireadout}
  \lean{PhyXMiniProblems.ProblemPhyXMini0821.nonnegativeSIReadout}
  Read a nonnegative dimensionful quantity in coherent SI units
\end{definition}
\begin{proof}
  This is the declared model definition of nonnegative SIReadout.
\end{proof}
\begin{definition}[mass In Kilograms]
  \label{def:physics:phyx-mini-0821:phyxminiproblems-problemphyxmini0821-massinkilograms}
  \lean{PhyXMiniProblems.ProblemPhyXMini0821.massInKilograms}
  \uses{def:physics:phyx-mini-0821:phyxminiproblems-problemphyxmini0821-massquantity, def:physics:phyx-mini-0821:phyxminiproblems-problemphyxmini0821-nonnegativesireadout}
  Kilogram readout of a mass
\end{definition}
\begin{proof}
  This is the declared model definition of mass In Kilograms.
\end{proof}
\begin{definition}[length In Meters]
  \label{def:physics:phyx-mini-0821:phyxminiproblems-problemphyxmini0821-lengthinmeters}
  \lean{PhyXMiniProblems.ProblemPhyXMini0821.lengthInMeters}
  \uses{def:physics:phyx-mini-0821:phyxminiproblems-problemphyxmini0821-lengthquantity, def:physics:phyx-mini-0821:phyxminiproblems-problemphyxmini0821-nonnegativesireadout}
  Metre readout of a length
\end{definition}
\begin{proof}
  This is the declared model definition of length In Meters.
\end{proof}
\begin{definition}[speed In Meters Per Second]
  \label{def:physics:phyx-mini-0821:phyxminiproblems-problemphyxmini0821-speedinmeterspersecond}
  \lean{PhyXMiniProblems.ProblemPhyXMini0821.speedInMetersPerSecond}
  \uses{def:physics:phyx-mini-0821:phyxminiproblems-problemphyxmini0821-speedquantity, def:physics:phyx-mini-0821:phyxminiproblems-problemphyxmini0821-nonnegativesireadout}
  Metre-per-second readout of a linear speed
\end{definition}
\begin{proof}
  This is the declared model definition of speed In Meters Per Second.
\end{proof}
\begin{definition}[angular Speed In Radians Per Second]
  \label{def:physics:phyx-mini-0821:phyxminiproblems-problemphyxmini0821-angularspeedinradianspersecond}
  \lean{PhyXMiniProblems.ProblemPhyXMini0821.angularSpeedInRadiansPerSecond}
  \uses{def:physics:phyx-mini-0821:phyxminiproblems-problemphyxmini0821-angularspeedquantity, def:physics:phyx-mini-0821:phyxminiproblems-problemphyxmini0821-nonnegativesireadout}
  Radian-per-second readout of an angular-speed magnitude
\end{definition}
\begin{proof}
  This is the declared model definition of angular Speed In Radians Per Second.
\end{proof}
\begin{definition}[acceleration In Meters Per Second Squared]
  \label{def:physics:phyx-mini-0821:phyxminiproblems-problemphyxmini0821-accelerationinmeterspersecondsquared}
  \lean{PhyXMiniProblems.ProblemPhyXMini0821.accelerationInMetersPerSecondSquared}
  \uses{def:physics:phyx-mini-0821:phyxminiproblems-problemphyxmini0821-accelerationquantity, def:physics:phyx-mini-0821:phyxminiproblems-problemphyxmini0821-nonnegativesireadout}
  Metre-per-second-squared readout of an acceleration magnitude
\end{definition}
\begin{proof}
  This is the declared model definition of acceleration In Meters Per Second Squared.
\end{proof}
\begin{definition}[moment Of Inertia In Kilogram Meters Squared]
  \label{def:physics:phyx-mini-0821:phyxminiproblems-problemphyxmini0821-momentofinertiainkilogrammeterssquared}
  \lean{PhyXMiniProblems.ProblemPhyXMini0821.momentOfInertiaInKilogramMetersSquared}
  \uses{def:physics:phyx-mini-0821:phyxminiproblems-problemphyxmini0821-momentofinertiaquantity, def:physics:phyx-mini-0821:phyxminiproblems-problemphyxmini0821-nonnegativesireadout}
  kg m² readout of a scalar moment of inertia
\end{definition}
\begin{proof}
  This is the declared model definition of moment Of Inertia In Kilogram Meters Squared.
\end{proof}
\begin{definition}[energy In Joules]
  \label{def:physics:phyx-mini-0821:phyxminiproblems-problemphyxmini0821-energyinjoules}
  \lean{PhyXMiniProblems.ProblemPhyXMini0821.energyInJoules}
  \uses{def:physics:phyx-mini-0821:phyxminiproblems-problemphyxmini0821-energyquantity}
  Joule readout of an energy, calibrated by Physlib's joule
\end{definition}
\begin{proof}
  This is the declared model definition of energy In Joules.
\end{proof}
\begin{definition}[Figure Element]
  \label{def:physics:phyx-mini-0821:phyxminiproblems-problemphyxmini0821-figureelement}
  \lean{PhyXMiniProblems.ProblemPhyXMini0821.FigureElement}
  Individually recognizable elements of the supplied bitmap
\end{definition}
\begin{proof}
  This is the declared model definition of Figure Element.
\end{proof}
\begin{definition}[Figure Label]
  \label{def:physics:phyx-mini-0821:phyxminiproblems-problemphyxmini0821-figurelabel}
  \lean{PhyXMiniProblems.ProblemPhyXMini0821.FigureLabel}
  Literal symbolic labels visible in the supplied bitmap
\end{definition}
\begin{proof}
  This is the declared model definition of Figure Label.
\end{proof}
\begin{definition}[Cylinder Side]
  \label{def:physics:phyx-mini-0821:phyxminiproblems-problemphyxmini0821-cylinderside}
  \lean{PhyXMiniProblems.ProblemPhyXMini0821.CylinderSide}
  Side of the cylinder at which the pictured straight cable is tangent
\end{definition}
\begin{proof}
  This is the declared model definition of Cylinder Side.
\end{proof}
\begin{definition}[Cable Segment Orientation]
  \label{def:physics:phyx-mini-0821:phyxminiproblems-problemphyxmini0821-cablesegmentorientation}
  \lean{PhyXMiniProblems.ProblemPhyXMini0821.CableSegmentOrientation}
  Orientation of the straight cable segment drawn between cylinder and block
\end{definition}
\begin{proof}
  This is the declared model definition of Cable Segment Orientation.
\end{proof}
\begin{definition}[Cylinder Cable Figure]
  \label{def:physics:phyx-mini-0821:phyxminiproblems-problemphyxmini0821-cylindercablefigure}
  \lean{PhyXMiniProblems.ProblemPhyXMini0821.CylinderCableFigure}
  \uses{def:physics:phyx-mini-0821:phyxminiproblems-problemphyxmini0821-figureelement, def:physics:phyx-mini-0821:phyxminiproblems-problemphyxmini0821-figurelabel, def:physics:phyx-mini-0821:phyxminiproblems-problemphyxmini0821-cylinderside, def:physics:phyx-mini-0821:phyxminiproblems-problemphyxmini0821-cablesegmentorientation}
  Qualitative placement information transcribed from primary image 821.png
\end{definition}
\begin{proof}
  This is the declared model definition of Cylinder Cable Figure.
\end{proof}
\begin{definition}[Cylinder Mass Model]
  \label{def:physics:phyx-mini-0821:phyxminiproblems-problemphyxmini0821-cylindermassmodel}
  \lean{PhyXMiniProblems.ProblemPhyXMini0821.CylinderMassModel}
  The rigid body's mass-distribution idealization
\end{definition}
\begin{proof}
  This is the declared model definition of Cylinder Mass Model.
\end{proof}
\begin{definition}[Axis Orientation]
  \label{def:physics:phyx-mini-0821:phyxminiproblems-problemphyxmini0821-axisorientation}
  \lean{PhyXMiniProblems.ProblemPhyXMini0821.AxisOrientation}
  Orientation of the cylinder's rotation axis
\end{definition}
\begin{proof}
  This is the declared model definition of Axis Orientation.
\end{proof}
\begin{definition}[Axis Motion]
  \label{def:physics:phyx-mini-0821:phyxminiproblems-problemphyxmini0821-axismotion}
  \lean{PhyXMiniProblems.ProblemPhyXMini0821.AxisMotion}
  Motion of the axis relative to the laboratory
\end{definition}
\begin{proof}
  This is the declared model definition of Axis Motion.
\end{proof}
\begin{definition}[Axle Friction Model]
  \label{def:physics:phyx-mini-0821:phyxminiproblems-problemphyxmini0821-axlefrictionmodel}
  \lean{PhyXMiniProblems.ProblemPhyXMini0821.AxleFrictionModel}
  Bearing-friction idealization used by the problem
\end{definition}
\begin{proof}
  This is the declared model definition of Axle Friction Model.
\end{proof}
\begin{definition}[Cable Mass Model]
  \label{def:physics:phyx-mini-0821:phyxminiproblems-problemphyxmini0821-cablemassmodel}
  \lean{PhyXMiniProblems.ProblemPhyXMini0821.CableMassModel}
  Cable-mass idealization represented by the word "light"
\end{definition}
\begin{proof}
  This is the declared model definition of Cable Mass Model.
\end{proof}
\begin{definition}[Cable Stretch Model]
  \label{def:physics:phyx-mini-0821:phyxminiproblems-problemphyxmini0821-cablestretchmodel}
  \lean{PhyXMiniProblems.ProblemPhyXMini0821.CableStretchModel}
  Cable-extension idealization represented by "nonstretching"
\end{definition}
\begin{proof}
  This is the declared model definition of Cable Stretch Model.
\end{proof}
\begin{definition}[Cable Contact Model]
  \label{def:physics:phyx-mini-0821:phyxminiproblems-problemphyxmini0821-cablecontactmodel}
  \lean{PhyXMiniProblems.ProblemPhyXMini0821.CableContactModel}
  Contact behavior between the cable and cylinder
\end{definition}
\begin{proof}
  This is the declared model definition of Cable Contact Model.
\end{proof}
\begin{definition}[Cable Attachment]
  \label{def:physics:phyx-mini-0821:phyxminiproblems-problemphyxmini0821-cableattachment}
  \lean{PhyXMiniProblems.ProblemPhyXMini0821.CableAttachment}
  Attachment of the cable's free end
\end{definition}
\begin{proof}
  This is the declared model definition of Cable Attachment.
\end{proof}
\begin{definition}[Cylinder Cable Setup]
  \label{def:physics:phyx-mini-0821:phyxminiproblems-problemphyxmini0821-cylindercablesetup}
  \lean{PhyXMiniProblems.ProblemPhyXMini0821.CylinderCableSetup}
  \uses{def:physics:phyx-mini-0821:phyxminiproblems-problemphyxmini0821-massquantity, def:physics:phyx-mini-0821:phyxminiproblems-problemphyxmini0821-lengthquantity, def:physics:phyx-mini-0821:phyxminiproblems-problemphyxmini0821-speedquantity, def:physics:phyx-mini-0821:phyxminiproblems-problemphyxmini0821-angularspeedquantity, def:physics:phyx-mini-0821:phyxminiproblems-problemphyxmini0821-accelerationquantity, def:physics:phyx-mini-0821:phyxminiproblems-problemphyxmini0821-momentofinertiaquantity, def:physics:phyx-mini-0821:phyxminiproblems-problemphyxmini0821-energyquantity, def:physics:phyx-mini-0821:phyxminiproblems-problemphyxmini0821-cylindercablefigure, def:physics:phyx-mini-0821:phyxminiproblems-problemphyxmini0821-cylindermassmodel, def:physics:phyx-mini-0821:phyxminiproblems-problemphyxmini0821-axisorientation, def:physics:phyx-mini-0821:phyxminiproblems-problemphyxmini0821-axismotion, def:physics:phyx-mini-0821:phyxminiproblems-problemphyxmini0821-axlefrictionmodel, def:physics:phyx-mini-0821:phyxminiproblems-problemphyxmini0821-cablemassmodel, def:physics:phyx-mini-0821:phyxminiproblems-problemphyxmini0821-cablestretchmodel, def:physics:phyx-mini-0821:phyxminiproblems-problemphyxmini0821-cablecontactmodel, def:physics:phyx-mini-0821:phyxminiproblems-problemphyxmini0821-cableattachment}
  Independent quantities and qualitative data of the cylinder-cable-block system. In particular, the impact angular speed is an unconstrained physical quantity here and is not defined from an answer choice
\end{definition}
\begin{proof}
  This is the declared model definition of Cylinder Cable Setup.
\end{proof}
\begin{definition}[Matches Primary Figure]
  \label{def:physics:phyx-mini-0821:phyxminiproblems-problemphyxmini0821-matchesprimaryfigure}
  \lean{PhyXMiniProblems.ProblemPhyXMini0821.MatchesPrimaryFigure}
  \uses{def:physics:phyx-mini-0821:phyxminiproblems-problemphyxmini0821-figureelement, def:physics:phyx-mini-0821:phyxminiproblems-problemphyxmini0821-figurelabel, def:physics:phyx-mini-0821:phyxminiproblems-problemphyxmini0821-cylindercablesetup}
  Literal graphical evidence from the primary raster
\end{definition}
\begin{proof}
  This is the declared model definition of Matches Primary Figure.
\end{proof}
\begin{definition}[Matches Problem Description]
  \label{def:physics:phyx-mini-0821:phyxminiproblems-problemphyxmini0821-matchesproblemdescription}
  \lean{PhyXMiniProblems.ProblemPhyXMini0821.MatchesProblemDescription}
  \uses{def:physics:phyx-mini-0821:phyxminiproblems-problemphyxmini0821-speedinmeterspersecond, def:physics:phyx-mini-0821:phyxminiproblems-problemphyxmini0821-angularspeedinradianspersecond, def:physics:phyx-mini-0821:phyxminiproblems-problemphyxmini0821-cylindercablesetup}
  Qualitative assumptions and initial/impact conditions stated in the prose
\end{definition}
\begin{proof}
  This is the declared model definition of Matches Problem Description.
\end{proof}
\begin{definition}[Has Physical Parameters]
  \label{def:physics:phyx-mini-0821:phyxminiproblems-problemphyxmini0821-hasphysicalparameters}
  \lean{PhyXMiniProblems.ProblemPhyXMini0821.HasPhysicalParameters}
  \uses{def:physics:phyx-mini-0821:phyxminiproblems-problemphyxmini0821-massinkilograms, def:physics:phyx-mini-0821:phyxminiproblems-problemphyxmini0821-lengthinmeters, def:physics:phyx-mini-0821:phyxminiproblems-problemphyxmini0821-accelerationinmeterspersecondsquared, def:physics:phyx-mini-0821:phyxminiproblems-problemphyxmini0821-cylindercablesetup}
  Positivity and nondegeneracy of the physical parameters
\end{definition}
\begin{proof}
  This is the declared model definition of Has Physical Parameters.
\end{proof}
\begin{definition}[Satisfies Cylinder Cable Laws]
  \label{def:physics:phyx-mini-0821:phyxminiproblems-problemphyxmini0821-satisfiescylindercablelaws}
  \lean{PhyXMiniProblems.ProblemPhyXMini0821.SatisfiesCylinderCableLaws}
  \uses{def:physics:phyx-mini-0821:phyxminiproblems-problemphyxmini0821-massinkilograms, def:physics:phyx-mini-0821:phyxminiproblems-problemphyxmini0821-lengthinmeters, def:physics:phyx-mini-0821:phyxminiproblems-problemphyxmini0821-speedinmeterspersecond, def:physics:phyx-mini-0821:phyxminiproblems-problemphyxmini0821-angularspeedinradianspersecond, def:physics:phyx-mini-0821:phyxminiproblems-problemphyxmini0821-accelerationinmeterspersecondsquared, def:physics:phyx-mini-0821:phyxminiproblems-problemphyxmini0821-momentofinertiainkilogrammeterssquared, def:physics:phyx-mini-0821:phyxminiproblems-problemphyxmini0821-energyinjoules, def:physics:phyx-mini-0821:phyxminiproblems-problemphyxmini0821-cylindercablesetup}
  Governing mechanics for the ideal system. The no-slip condition is kept in its physical rim-speed form through an independently represented cable speed; the requested solved relation ω = v / R does not occur in this structure
\end{definition}
\begin{proof}
  This is the declared model definition of Satisfies Cylinder Cable Laws.
\end{proof}
\begin{lemma}[block Impact Speed eq radius mul angular Speed]
  \label{lem:physics:phyx-mini-0821:phyxminiproblems-problemphyxmini0821-blockimpactspeed-eq-radius-mul-angularspeed}
  \lean{PhyXMiniProblems.ProblemPhyXMini0821.blockImpactSpeed_eq_radius_mul_angularSpeed}
  \uses{def:physics:phyx-mini-0821:phyxminiproblems-problemphyxmini0821-lengthinmeters, def:physics:phyx-mini-0821:phyxminiproblems-problemphyxmini0821-speedinmeterspersecond, def:physics:phyx-mini-0821:phyxminiproblems-problemphyxmini0821-angularspeedinradianspersecond, def:physics:phyx-mini-0821:phyxminiproblems-problemphyxmini0821-cylindercablesetup, def:physics:phyx-mini-0821:phyxminiproblems-problemphyxmini0821-satisfiescylindercablelaws}
  The cable constraints combine to identify the block speed with rim speed
\end{lemma}
\begin{proof}
  The conclusion follows from the governing relations and figure readouts named in the statement of block Impact Speed eq radius mul angular Speed.
\end{proof}
\begin{definition}[Answer Choice]
  \label{def:physics:phyx-mini-0821:phyxminiproblems-problemphyxmini0821-answerchoice}
  \lean{PhyXMiniProblems.ProblemPhyXMini0821.AnswerChoice}
  Labels of the four formulas printed as answer choices
\end{definition}
\begin{proof}
  This is the declared model definition of Answer Choice.
\end{proof}
\begin{definition}[recorded Answer Choice]
  \label{def:physics:phyx-mini-0821:phyxminiproblems-problemphyxmini0821-recordedanswerchoice}
  \lean{PhyXMiniProblems.ProblemPhyXMini0821.recordedAnswerChoice}
  \uses{def:physics:phyx-mini-0821:phyxminiproblems-problemphyxmini0821-answerchoice}
  Dataset metadata recording the supplied answer label
\end{definition}
\begin{proof}
  This is the declared model definition of recorded Answer Choice.
\end{proof}
% --- Archon declaration topology end ---
% --- Archon physics formalization source end ---
